
\section{Experimental Setup}

\subsection{Dataset}

\hspace*{0.3in} The SYNAPSE platform is evaluated using a combination of real-world telemetry streams and controlled synthetic fault scenarios. Real operational metrics are collected through a Prometheus scraping pipeline integrated with a ten-node Kubernetes microservice cluster. Each node exposes five core time-series signals — \textbf{latency} (p95, in milliseconds), \textbf{error\_rate} (proportion of 5xx responses), \textbf{cpu\_usage} (percentage of allocated cores), \textbf{memory\_usage} (percentage of container memory limit), and \textbf{request\_rate} (requests per second) — sampled at fifteen-second intervals and buffered through an in-process ingestion queue before entering the AI pipeline.\\\\

\hspace*{0.3in} Controlled fault injection is performed by the internal \textbf{MicroserviceSimulator} module, which programmatically introduces five canonical fault archetypes into the live metric stream: \textbf{latency\_spike} (sudden p95 latency surge beyond a configurable threshold), \textbf{error\_surge} (rapid increase in HTTP error ratio), \textbf{cpu\_throttle} (CPU contention causing service degradation), \textbf{memory\_leak} (monotonically growing heap consumption), and \textbf{network\_partition} (dropped inter-service connectivity resulting in timeout cascades). The simulated dependency graph comprises ten representative microservice nodes: \textit{API Gateway}, \textit{Auth Service}, \textit{User Service}, \textit{Product Service}, \textit{Order Service}, \textit{Payment Service}, \textit{Notification Service}, \textit{Analytics Service}, \textit{Cache Layer}, and \textit{Database}. These nodes are arranged according to a realistic directed call graph where upstream failures propagate downstream, enabling genuine cascading fault scenarios.\\\\

\hspace*{0.3in} The evaluation corpus consists of \textbf{500 normal baseline scenarios} used for Graph Neural Network (GNN) pre-training, each representing fifteen-minute windows of healthy cluster operation, and \textbf{200 annotated fault injection cases} used for quantitative evaluation. Every fault case carries a ground-truth label specifying the root-cause service, the fault archetype, and the propagation path, enabling precision/recall-based assessment of the full RCA pipeline. Key features extracted from this dataset fall into three categories:

\begin{itemize}
    \item \textbf{Metrics Features:} Per-node z-score normalized vectors of the five Prometheus signals, computed over a rolling window of twenty data points, yielding a $10 \times 5$ node-feature matrix fed to the VAE-GNN encoder at each inference step.
    \item \textbf{Graph Features:} Directed adjacency matrix encoding service dependencies, augmented with dynamic edge weights derived from real-time inter-service call latency and error rates, enabling the GNN message-passing layers to propagate anomaly signals along actual dependency paths.
    \item \textbf{Incident Ground Truth:} Annotated incident records stored in the \texttt{Incident} database table, each containing the injected fault type, the causal root-cause service, the observed severity level (\texttt{L1\_SIMPLE} through \texttt{L4\_DESIGN\_FLAW}), and a timestamped propagation trail used to validate the DECI causal graph output.
\end{itemize}

\subsection{Performance Parameters}

\hspace*{0.3in} Six primary performance parameters govern the quantitative evaluation of SYNAPSE. These parameters collectively assess the correctness, speed, resource footprint, and long-term reliability of the platform across the full ten-stage AI pipeline.

\begin{enumerate}[label=\textbf{\Alph*.},leftmargin=1.2cm]

    \item \textbf{RCA Accuracy} refers to the fraction of evaluated fault scenarios for which the platform correctly identifies the ground-truth root-cause service as the top-ranked candidate in the \texttt{RCAResult} output. A prediction is considered correct when the identified root-cause service and the fault archetype both match the annotated ground truth. RCA Accuracy is measured by comparing SYNAPSE's top-1 causal attribution against the 200-case annotated evaluation set. The design target for this metric is \textbf{85\% or higher}, reflecting the practical requirement that an AIOps tool must outperform heuristic rule-based baselines to justify autonomous deployment in production environments.

    \item \textbf{Pipeline Latency} quantifies the end-to-end wall-clock time elapsed from the moment a metrics payload enters the ingestion endpoint to the moment a finalized \texttt{RCAResult} record, complete with LLM-generated reasoning and recovery action recommendations, is persisted to the database. The measurement spans all five major computational stages: (i) metrics ingestion and feature normalization, (ii) VAE-GNN anomaly detection and anomaly score generation, (iii) DECI causal discovery and DAG construction, (iv) GPT-4o-mini LLM reasoning synthesis, and (v) recovery action selection and database persistence. The target is a total end-to-end latency of \textbf{less than 10 seconds} under nominal load, ensuring that SYNAPSE can provide actionable RCA outputs well within the incident response window expected by on-call engineering teams.

    \item \textbf{GNN Anomaly Detection F1-Score} evaluates the binary classification performance of the VAE-GNN anomaly detection stage in isolation. Precision measures the fraction of flagged anomalous nodes that are genuinely faulty; recall measures the fraction of all genuinely faulty nodes that are successfully flagged. The F1-Score, defined as the harmonic mean of precision and recall, serves as the primary aggregate indicator of anomaly detector quality. This metric is computed per fault injection scenario across the 200-case evaluation corpus, and the macro-averaged F1 across all scenarios and fault types is reported as the headline GNN performance figure. The target F1-Score is \textbf{0.88 or above}.

    \item \textbf{Resource Utilization} tracks the computational footprint of the SYNAPSE backend during sustained operation. GPU utilization (where available, via CUDA) is monitored during GNN batch inference; CPU utilization is measured across all pipeline stages using process-level sampling; and heap memory consumption is tracked to detect leaks under long-duration runs. Resource utilization is measured using Python's \texttt{psutil} library and, when applicable, \texttt{nvidia-smi} for GPU memory reporting. The target is that the platform remains operable on a single server instance equipped with 8 CPU cores, 16 GB RAM, and an optional mid-range GPU (e.g., NVIDIA T4), without exceeding 70\% sustained CPU utilization or 80\% GPU memory occupancy during peak inference bursts.

    \item \textbf{Continual Learning Performance} assesses the stability-plasticity balance of the continual learning module, which combines Elastic Weight Consolidation (EWC) and Model-Agnostic Meta-Learning (MAML) with experience replay. Three sub-metrics are reported: (i) \textbf{Forgetting Rate} — the degradation in RCA accuracy on previously encountered fault patterns after training on new fault types, measured as the absolute accuracy drop; (ii) \textbf{MAML Adaptation Steps} — the number of gradient steps required for the model to achieve acceptable performance on a novel fault scenario in the few-shot adaptation phase, with a target of five steps or fewer; and (iii) \textbf{EWC Lambda Effectiveness} — whether the chosen regularization coefficient $\lambda$ successfully prevents catastrophic forgetting while still permitting adequate plasticity on new patterns, validated by comparing accuracy before and after EWC-regularized updates.

    \item \textbf{Model Reproducibility} verifies that the SYNAPSE pipeline produces identical \texttt{RCAResult} outputs — including root-cause ranking, severity classification, and recovery action selection — when presented with the same input metrics payload under identical model weights and random seeds. Reproducibility is validated by replaying ten randomly selected fault scenarios five times each and confirming that all outputs match across runs. This parameter ensures that SYNAPSE's recommendations are deterministic and auditable, a critical property for post-incident forensic analysis and regulatory compliance in production deployments.

\end{enumerate}

\subsection{Efficiency Issues}

\hspace*{0.3in} Several computational bottlenecks were identified during the design and integration testing phases of SYNAPSE. The following four issues represent the most significant efficiency challenges encountered, along with the mitigation strategies adopted in the final implementation.

\begin{enumerate}[label=\textbf{\Alph*.},leftmargin=1.2cm]

    \item \textbf{GNN Inference Overhead:} The VAE-GNN model, when executed on CPU for each incoming incident individually, introduced non-trivial inference latency due to the repeated construction and forward-passing of the graph tensor for every request. Under concurrent incident load, this overhead compounded and risked breaching the ten-second pipeline budget. This issue was mitigated by introducing \textbf{batch processing} within the GNN inference stage, grouping concurrent anomaly detection requests into mini-batches before invoking the PyTorch Geometric (PyG) forward pass. Additionally, the graph tensor construction was refactored to use PyG's optimized \texttt{Data} and \texttt{Batch} primitives, eliminating redundant memory allocations and leveraging PyG's built-in sparse adjacency optimizations for the ten-node dependency graph.

    \item \textbf{DECI Causal Discovery Latency:} The DECI (Differentiable Causal Discovery from Interventional Data) causal graph construction stage is computationally intensive, as it involves iterative optimization of a structural equation model over the service dependency graph. Applying DECI to the full ten-node graph for every incident introduced latency spikes that threatened real-time responsiveness. The adopted mitigation strategy restricts DECI inference exclusively to the \textbf{anomalous subset} of nodes flagged by the preceding GNN stage, typically comprising two to four nodes per incident, which reduces the effective graph size and computation time substantially. A configurable \textbf{timeout fallback} mechanism was also implemented, such that if DECI does not converge within a predefined budget (default: three seconds), the pipeline falls back to a correlation-based heuristic causal ranking derived from the RAG-constrained $W_0$ prior.

    \item \textbf{LLM API Latency:} Integration with the GPT-4o-mini API for natural-language RCA reasoning synthesis introduced variable network round-trip latency, which in adversarial network conditions could constitute the dominant contributor to total pipeline time. Three complementary mitigations were adopted: (i) a \textbf{mock mode} for development and testing environments that returns a pre-cached LLM response without making any API call, enabling full pipeline testing without incurring API costs or latency; (ii) \textbf{asynchronous API invocation} using Python's \texttt{asyncio} and \texttt{httpx} libraries, ensuring that the LLM call does not block other pipeline stages; and (iii) \textbf{response caching} keyed on a hash of the LLM prompt context, so that semantically identical incidents do not trigger redundant API calls.

    \item \textbf{Database I/O Contention:} Synchronous database writes to the \texttt{RCAResult}, \texttt{ServiceMetric}, \texttt{RecoveryActionLog}, and \texttt{PatternHistory} tables introduced blocking delays during peak incident ingestion, particularly when multiple concurrent RCA pipelines attempted simultaneous persistence. This issue was resolved by using \textbf{SQLAlchemy's asynchronous session} interface (\texttt{AsyncSession}) throughout the persistence layer, decoupling database I/O from the pipeline's critical path. Additionally, \textbf{batch write} operations were introduced for \texttt{ServiceMetric} records, consolidating multiple per-node metric inserts into a single bulk \texttt{INSERT} statement, significantly reducing the number of round-trips to the SQLite/PostgreSQL backend under high-throughput scenarios.

\end{enumerate}

% ─────────────────────────────────────────────────────────────────────────────
\section{Software Testing}

\hspace*{0.3in} Software testing constitutes a cornerstone of the SYNAPSE development lifecycle, ensuring that each individual component, inter-module interface, and end-to-end pipeline pathway functions correctly under both nominal operating conditions and adversarial edge cases. Given that SYNAPSE operates as an automated decision-making system that directly influences infrastructure recovery actions in production Kubernetes environments, defects in any layer of the stack — from metric ingestion to LLM reasoning to recovery command generation — carry the potential for incorrect automated remediation, making comprehensive test coverage a non-negotiable engineering requirement.\\\\

\hspace*{0.3in} The SYNAPSE test suite spans multiple testing disciplines: \textbf{unit testing} validates individual functions and model inference routines in isolation; \textbf{integration testing} verifies the correct data flow across adjacent pipeline stages and between the FastAPI backend and the SQLAlchemy persistence layer; \textbf{system testing} exercises the full ten-stage pipeline from raw metric ingestion to database-persisted \texttt{RCAResult}; and \textbf{performance testing} benchmarks latency and resource consumption under simulated concurrent load. All test cases are designed to produce deterministic, repeatable outcomes and are executed against a sandboxed environment populated with the synthetic \texttt{MicroserviceSimulator} dataset, ensuring that no live production infrastructure is affected during evaluation runs.

\subsection{Test Cases and Test Results}

\renewcommand{\arraystretch}{1.3}
\small
\begin{center}
\begin{longtable}{|c|>{\RaggedRight\arraybackslash}p{3.2cm}|>{\RaggedRight\arraybackslash}p{3.5cm}|>{\RaggedRight\arraybackslash}p{3.5cm}|>{\RaggedRight\arraybackslash}p{2.2cm}|c|}
\caption{Performance Evaluation Test Cases for SYNAPSE} \label{tab:testcases} \\
\hline
\textbf{TC No.} & \textbf{Test Case Description} & \textbf{Test Input} & \textbf{Expected Output} & \textbf{Actual Output} & \textbf{Status} \\
\hline
\endfirsthead
\multicolumn{6}{c}{\tablename\ \thetable{} -- \textit{Continued from previous page}} \\
\hline
\textbf{TC No.} & \textbf{Test Case Description} & \textbf{Test Input} & \textbf{Expected Output} & \textbf{Actual Output} & \textbf{Status} \\
\hline
\endhead
\hline \multicolumn{6}{r}{\textit{Continued on next page}} \\
\endfoot
\hline
\endlastfoot

TC-001 & Valid metrics ingestion via POST /api/v1/metrics/ingest & JSON payload with 10-node metrics, all fields populated & HTTP 200 OK, incident record created in DB & HTTP 200 OK, incident stored & \textbf{Pass} \\
\hline

TC-002 & Invalid metrics ingestion — missing required field & JSON payload with \texttt{service\_name} field omitted & HTTP 422 Unprocessable Entity with field error detail & HTTP 422 with Pydantic validation error & \textbf{Pass} \\
\hline

TC-003 & Invalid metrics ingestion — out-of-range value & Payload with \texttt{cpu\_usage} = 150.0 (exceeds 100\%) & HTTP 422 with range validation error message & HTTP 422 returned as expected & \textbf{Pass} \\
\hline

TC-004 & GNN anomaly detection on normal baseline scenario & 10-node feature matrix sampled from normal baseline & All anomaly scores below threshold; no anomalies flagged & Zero anomalous nodes detected & \textbf{Pass} \\
\hline

TC-005 & GNN anomaly detection on injected latency spike & 10-node feature matrix with latency\_spike on Payment Service & Payment Service node flagged as anomalous; score above 0.7 & Payment Service flagged, score 0.83 & \textbf{Pass} \\
\hline

TC-006 & GNN anomaly detection on memory leak scenario & Feature matrix with monotonically rising memory on DB node & DB node flagged; adjacent downstream nodes show elevated scores & DB node flagged, cascade scores elevated & \textbf{Pass} \\
\hline

TC-007 & Severity classification — L1 simple fault & Single-node anomaly with low error rate and moderate latency & Severity assigned as \texttt{L1\_SIMPLE} & \texttt{L1\_SIMPLE} returned & \textbf{Pass} \\
\hline

TC-008 & Severity classification — L2 cascading fault & Two-node anomaly chain (Order → Payment) with elevated error rates & Severity assigned as \texttt{L2\_CASCADING} & \texttt{L2\_CASCADING} returned & \textbf{Pass} \\
\hline

TC-009 & Severity classification — L3 systemic fault & Five-node anomaly with cluster-wide CPU throttling & Severity assigned as \texttt{L3\_SYSTEMIC} & \texttt{L3\_SYSTEMIC} returned & \textbf{Pass} \\
\hline

TC-010 & Severity classification — L4 design flaw & Repeated pattern across 10+ incidents matching known anti-pattern & Severity assigned as \texttt{L4\_DESIGN\_FLAW} & \texttt{L4\_DESIGN\_FLAW} returned & \textbf{Pass} \\
\hline

TC-011 & DECI causal discovery on two-node anomalous subgraph & Anomalous node features for Auth and User services & DAG edge Auth $\rightarrow$ User inferred; Auth ranked as root cause & Auth $\rightarrow$ User DAG produced & \textbf{Pass} \\
\hline

TC-012 & DECI causal discovery — timeout fallback activation & DECI configured with 0.1 s timeout on 6-node subgraph & Pipeline falls back to correlation-based heuristic causal ranking & Fallback heuristic ranking returned within 0.5 s & \textbf{Pass} \\
\hline

TC-013 & DECI output consistency with RAG-constrained $W_0$ prior & Known dependency structure injected as prior; anomaly on Order node & Causal graph respects prior edges; Order Service as root cause & Prior-consistent DAG, Order as root cause & \textbf{Pass} \\
\hline

TC-014 & LLM reasoning output — valid RCA summary generation & Structured RCA context JSON sent to GPT-4o-mini reasoning stage & Natural-language explanation with root cause, impact, and steps & Coherent explanation generated, root cause named & \textbf{Pass} \\
\hline

TC-015 & LLM reasoning — mock mode response & LLM stage configured in \texttt{MOCK\_LLM=true} mode & Pre-cached mock response returned; no real API call made & Mock response returned in $<$50 ms & \textbf{Pass} \\
\hline

TC-016 & LLM reasoning — response caching on repeated incident & Same RCA context submitted twice in succession & Second call returns cached response; API not invoked again & Cache hit confirmed, latency $<$10 ms on second call & \textbf{Pass} \\
\hline

TC-017 & Recovery action selection — pod restart for L1 fault & \texttt{L1\_SIMPLE} RCA result on Payment Service & \texttt{kubectl rollout restart deployment/payment-service} command generated & Correct kubectl command produced & \textbf{Pass} \\
\hline

TC-018 & Recovery action selection — HPA scale-out for CPU throttle & \texttt{L2\_CASCADING} fault with CPU throttle root cause & Horizontal Pod Autoscaler scale-out command generated for affected service & HPA scale-out command generated correctly & \textbf{Pass} \\
\hline

TC-019 & Recovery action selection — circuit breaker for network partition & Network partition fault on inter-service link & Circuit breaker activation recommendation generated & Circuit breaker kubectl patch command returned & \textbf{Pass} \\
\hline

TC-020 & Fault simulation — latency spike injection via API & POST /api/v1/simulate/fault with type=latency\_spike, target=order-service & Fault injected into metric stream; incident created within 5 s & Incident created in 3.2 s & \textbf{Pass} \\
\hline

TC-021 & Fault simulation — error surge injection & POST /api/v1/simulate/fault with type=error\_surge, target=auth-service & Elevated error\_rate metrics generated; GNN flags Auth node & Auth node flagged by GNN & \textbf{Pass} \\
\hline

TC-022 & Incident CRUD — create incident record & POST /api/v1/incidents with valid incident payload & Incident persisted to DB; HTTP 201 Created with new ID returned & HTTP 201, incident ID 42 returned & \textbf{Pass} \\
\hline

TC-023 & Incident CRUD — read incident by ID & GET /api/v1/incidents/42 & HTTP 200 with complete incident JSON including RCA result & HTTP 200, full incident JSON returned & \textbf{Pass} \\
\hline

TC-024 & Incident CRUD — update incident status & PATCH /api/v1/incidents/42 with \texttt{status=resolved} & HTTP 200, incident status updated to resolved in DB & Status updated, timestamp recorded & \textbf{Pass} \\
\hline

TC-025 & Incident CRUD — delete incident record & DELETE /api/v1/incidents/42 & HTTP 204 No Content; incident removed from DB & HTTP 204, record deleted & \textbf{Pass} \\
\hline

TC-026 & WebSocket live metrics stream & Client connects to \texttt{ws://host/ws/metrics/live} & Continuous JSON metric frames broadcast at 15 s intervals & Frames received at 14.8 s mean interval & \textbf{Pass} \\
\hline

TC-027 & WebSocket live RCA event push & Active RCA pipeline completes during WebSocket session & RCA result event pushed to all connected clients in real time & RCA event delivered within 200 ms of pipeline completion & \textbf{Pass} \\
\hline

TC-028 & Continual learning — EWC regularization prevents forgetting & Model trained on fault-type A then fine-tuned on fault-type B with EWC $\lambda=1000$ & Accuracy on fault-type A degrades by less than 5\% after B training & 3.1\% forgetting measured & \textbf{Pass} \\
\hline

TC-029 & Continual learning — MAML few-shot adaptation & 5 gradient steps on 3 novel fault examples & Model achieves $>$80\% accuracy on novel fault after adaptation & 84\% accuracy after 5 steps & \textbf{Pass} \\
\hline

TC-030 & Continual learning — experience replay buffer & Replay buffer sampled during EWC training cycle & Historical samples correctly retrieved and included in training batch & Replay samples included, loss balanced & \textbf{Pass} \\
\hline

TC-031 & Model training endpoint — trigger retraining & POST /api/v1/model/train with new labelled dataset & Training job launched; HTTP 202 Accepted with job ID returned & HTTP 202, job ID returned & \textbf{Pass} \\
\hline

TC-032 & Model status endpoint & GET /api/v1/model/status & HTTP 200 with current model version, last trained timestamp, and accuracy metrics & All fields present in response & \textbf{Pass} \\
\hline

TC-033 & Health check endpoint & GET /api/v1/health & HTTP 200 with \texttt{\{"status":"healthy"\}} and DB connectivity confirmation & HTTP 200, healthy status confirmed & \textbf{Pass} \\
\hline

TC-034 & API response time under concurrent load & 20 simultaneous POST /api/v1/metrics/ingest requests & All 20 requests return HTTP 200; mean response time $<$500 ms & Mean 387 ms, all 200 OK & \textbf{Pass} \\
\hline

TC-035 & Docker deployment — container startup and pipeline readiness & \texttt{docker-compose up} on clean environment & All services start; health check returns healthy within 30 s & Healthy in 22 s & \textbf{Pass} \\

\end{longtable}
\end{center}

% ─────────────────────────────────────────────────────────────────────────────
\section{Summary}

\hspace*{0.3in} The experimental evaluation of SYNAPSE demonstrates that the platform successfully fulfils its core design mandate of delivering automated, explainable, and actionable Root Cause Analysis for microservice failure incidents. Across the 200-case annotated evaluation corpus, SYNAPSE achieves a top-1 RCA accuracy of \textbf{89.5\%}, surpassing the 85\% design threshold by a meaningful margin and substantially outperforming baseline heuristic approaches that rely solely on threshold-based alerting. The VAE-GNN anomaly detection stage attains a macro-averaged F1-Score of \textbf{0.912} across all five fault archetypes, demonstrating that the graph-structured anomaly detector effectively discriminates faulty nodes from healthy ones even when fault signals propagate through multiple dependency hops.\\\\

\hspace*{0.3in} From an operational performance perspective, the end-to-end pipeline latency averages \textbf{6.8 seconds} under nominal single-incident load, comfortably within the ten-second budget. The dominant contributor to this latency is the GPT-4o-mini LLM reasoning stage (mean 3.1 s), which is partially offset by asynchronous invocation and response caching. The DECI causal discovery stage, restricted to the anomalous node subset, contributes an average of 1.4 seconds, while GNN inference, ingestion, and database persistence collectively account for the remaining 2.3 seconds. Resource utilization measurements confirm that the platform operates within its target footprint: sustained CPU utilization remains below 58\% on an 8-core host, and heap memory consumption stabilizes below 4.2 GB during continuous operation across a 72-hour soak test.\\\\

\hspace*{0.3in} The continual learning module demonstrates effective management of the stability-plasticity tradeoff. EWC regularization with $\lambda = 1000$ limits the forgetting rate on previously learned fault patterns to an average of \textbf{3.1\%} accuracy degradation after exposure to new fault types, well below the 5\% acceptable threshold. MAML few-shot adaptation consistently achieves over 80\% accuracy on novel fault archetypes within five gradient steps, confirming that the meta-learning component genuinely accelerates adaptation to previously unseen failure modes. Experience replay further reinforces retention, with replay-augmented training exhibiting consistently lower variance in accuracy across sequential task exposures compared to baseline fine-tuning without replay.\\\\

\hspace*{0.3in} The software testing campaign validates the reliability and correctness of every major functional boundary within SYNAPSE. All 35 test cases across the complete test suite — spanning metrics ingestion, GNN inference, severity classification, causal discovery, LLM reasoning, recovery action generation, WebSocket streaming, database CRUD, continual learning, and deployment — achieve a \textbf{Pass} status. API response time testing under 20 concurrent requests yields a mean latency of 387 milliseconds with zero error responses, confirming that the FastAPI asynchronous backend maintains acceptable responsiveness under realistic concurrent load. Docker-based deployment testing confirms that the complete multi-container stack achieves readiness within 22 seconds of a clean cold start.

\subsection{Result Analysis and Discussion}

\hspace*{0.3in} The achieved RCA accuracy of \textbf{89.5\%} represents a significant advancement over prior art in microservice fault diagnosis. Rule-based AIOps tools that rely on static threshold alerting typically achieve RCA accuracy in the range of 55–65\% on realistic cascading failure scenarios, as they cannot distinguish between the originating fault and downstream symptom manifestations. Graph-based approaches without causal reasoning have been reported in the literature at 72–78\% accuracy on comparable benchmark graphs. SYNAPSE's integration of VAE-GNN anomaly detection with DECI causal discovery and RAG-constrained priors bridges this gap, with the RAG prior contributing an estimated 4.2 percentage points of accuracy lift over an unconstrained DECI baseline, as measured by ablation experiments where the $W_0$ prior was set to a uniform uninformative matrix.\\\\

\hspace*{0.3in} The GNN F1-Score of \textbf{0.912} macro-averaged across fault types reveals notable variation across individual archetypes. The \textit{latency\_spike} archetype achieves the highest individual F1 of 0.943, benefiting from clear signal-to-noise separation in the latency feature dimension. The \textit{memory\_leak} archetype produces the lowest individual F1 of 0.871, attributable to the gradual onset characteristic of memory leaks, which makes early detection harder before the anomaly score crosses the decision boundary. The \textit{network\_partition} archetype yields F1 = 0.908, where the primary challenge is that the directly affected service may initially appear healthy in its own metrics while adjacent services exhibit error surges — a scenario where the GNN's message-passing mechanism, by aggregating neighbor features, proves essential for correct fault localization.\\\\

\hspace*{0.3in} Pipeline latency profiling across 200 evaluation runs reveals that the \textbf{99th percentile} end-to-end latency is \textbf{14.3 seconds}, driven primarily by LLM API tail latency under network congestion. The response caching mechanism reduces LLM calls by \textbf{31\%} across the evaluation corpus, as semantically similar incidents — particularly repeated \textit{latency\_spike} or \textit{error\_surge} scenarios on the same service — produce cache-hit prompts. The DECI timeout fallback activates in approximately \textbf{4\%} of cases, always on complex multi-node anomaly subgraphs of six or more nodes, and the heuristic fallback produces correct root-cause attribution in 71\% of these cases — lower than the full DECI accuracy of 89.5\% but still substantially better than random attribution. These figures collectively suggest that the timeout fallback, while a latency-driven compromise, preserves acceptable RCA quality in the cases where it is invoked.\\\\

\hspace*{0.3in} LLM reasoning quality was assessed through a blind evaluation in which three domain experts rated the generated natural-language RCA summaries on a five-point Likert scale across three dimensions: technical accuracy, actionability of recommendations, and clarity of explanation. The mean scores were \textbf{4.3/5.0} for technical accuracy, \textbf{4.1/5.0} for actionability, and \textbf{4.6/5.0} for clarity. Inter-rater agreement, measured by Fleiss' kappa, was $\kappa = 0.71$, indicating substantial agreement. The LLM-generated recovery actions achieved a Jaccard similarity of \textbf{0.78} against expert-curated reference remediation steps, confirming that the prompt engineering and RCA context structuring approach produces recommendations closely aligned with what experienced SRE practitioners would prescribe for the same incident type.\\\\

\hspace*{0.3in} Scalability experiments conducted by incrementally scaling the simulated microservice graph from 10 to 50 nodes reveal that GNN inference time grows sub-linearly with graph size due to the sparse adjacency structure and PyG batch optimization, increasing from a mean of 0.41 seconds at 10 nodes to 1.87 seconds at 50 nodes — a \textbf{4.6$\times$} latency increase for a \textbf{5$\times$} node count increase. DECI causal discovery, restricted to the anomalous subset, shows near-constant latency growth as long as the anomalous subset size remains small relative to the full graph, validating the design decision to apply DECI selectively rather than globally. Database write throughput sustains \textbf{320 RCA pipeline completions per minute} under batch write mode without degradation, confirming that the asynchronous SQLAlchemy persistence layer scales adequately for enterprise-scale microservice clusters with frequent incident generation. These results collectively establish SYNAPSE as a practically viable AIOps solution capable of deployment in real-world, large-scale Kubernetes production environments.
